% GENERATED FILE. Edit canonical lessons, then run npm run generate:books.
% canonical-source-hash: fnv1a64:f99a19d2cee7fd87
% canonical-lessons: FR-C372-risque, FR-C372-hasardeux, FR-C372-improbable, FR-C372-probabilite, FR-C372-quitte-a

\chapter{Prendre un risque}
\label{ch:fr-prendre-un-risque}
\hlchaptermodality{372}

\begin{chapteropening}
\textbf{By the end of this chapter:} \emph{I can weigh a risk: say something is chancy or unlikely, talk about the probability, and accept a cost.}

Complete the last lesson of chapter 372: I can weigh a risk: say something is chancy or unlikely, talk about the probability, and accept a cost.
\end{chapteropening}

\section[le risque]{le risque — the risk}

\label{lesson:FR-C372-risque}



\begin{quote}
Before the new one: say the French for final, then the French for to schedule.
\end{quote}

\subsection*{You'll want to know: le risque}
\textbf{le risque} — \textquotedblleft{}the risk\textquotedblright{}.

\begin{quote}
Il y a un risque de pluie demain.
\end{quote}

\begin{quote}
\emph{There is a chance of rain tomorrow.}
\end{quote}

\emph{Prendre un risque} — to take a risk; the verb is \emph{risquer}.

\begin{cousinweb}
\emph{risque} comes from Old Italian \emph{risco} \textquotedblleft{}danger\textquotedblright{} (EtymDB). The further origin is not settled.
\end{cousinweb}

\begin{grammarlens}[title={What you've built}]
One more word for telling, explaining and describing at B1.
\end{grammarlens}

\subsection*{Your turn}
\begin{itemize}
\raggedright
  \item \emph{Say it:} \emph{le risque}
  \item \emph{Say it:} \emph{le risque}, once more
  \item \emph{Say it:} a sentence of your own with \emph{le risque}
  \item \emph{Recall:} say the French for final, then the French for to schedule, then say \emph{le risque} again
\end{itemize}

\begin{tcolorbox}[breakable,colback=teal!4,colframe=teal!35!black,title={Before you move on}]
What does le risque mean? (\textquotedblleft{}The risk\textquotedblright{}.) Say it once more.
\end{tcolorbox}
\section[hasardeux]{hasardeux — risky, chancy}

\label{lesson:FR-C372-hasardeux}



\begin{quote}
Before the new one: say the French for to schedule, then the French for the risk.
\end{quote}

\subsection*{You'll want to know: hasardeux}
\textbf{hasardeux} — \textquotedblleft{}risky, chancy\textquotedblright{}.

\begin{quote}
C'est un projet hasardeux, mais il veut essayer.
\end{quote}

\begin{quote}
\emph{It is a risky project, but he wants to try.}
\end{quote}

Feminine \emph{hasardeuse}. Formal; in speech, \emph{risqué}.

\begin{cousinweb}
\emph{hasardeux} is built from an old French word for chance and the suffix \emph{-eux}. That word is usually traced, through Spanish, to Arabic \emph{az-zahr} \textquotedblleft{}the die\textquotedblright{}; the details are not settled.
\end{cousinweb}

\begin{grammarlens}[title={What you've built}]
One more word for telling, explaining and describing at B1.
\end{grammarlens}

\subsection*{Your turn}
\begin{itemize}
\raggedright
  \item \emph{Say it:} \emph{hasardeux}
  \item \emph{Say it:} \emph{hasardeux}, once more
  \item \emph{Say it:} a sentence of your own with \emph{hasardeux}
  \item \emph{Recall:} say the French for to schedule, then the French for the risk, then say \emph{hasardeux} again
\end{itemize}

\begin{tcolorbox}[breakable,colback=teal!4,colframe=teal!35!black,title={Before you move on}]
What does hasardeux mean? (\textquotedblleft{}Risky, chancy\textquotedblright{}.) Say it once more.
\end{tcolorbox}
\section[improbable]{improbable — improbable, unlikely}

\label{lesson:FR-C372-improbable}



\begin{quote}
Before the new one: say the French for the risk, then the French for risky.
\end{quote}

\subsection*{You'll want to know: improbable}
\textbf{improbable} — \textquotedblleft{}improbable, unlikely\textquotedblright{}.

\begin{quote}
Il est improbable qu'il vienne ce soir.
\end{quote}

\begin{quote}
\emph{It is unlikely that he will come tonight.}
\end{quote}

\emph{Il est improbable que} takes the subjunctive.

\begin{cousinweb}
\emph{improbable} is a learned borrowing from Latin \emph{improbābilis} \textquotedblleft{}not to be approved\textquotedblright{} (Lewis \& Short), built from \emph{in-} \textquotedblleft{}not\textquotedblright{} and \emph{probābilis}, from \emph{probāre} \textquotedblleft{}to test\textquotedblright{}.
\end{cousinweb}

\begin{grammarlens}[title={What you've built}]
One more word for telling, explaining and describing at B1.
\end{grammarlens}

\subsection*{Your turn}
\begin{itemize}
\raggedright
  \item \emph{Say it:} \emph{improbable}
  \item \emph{Say it:} \emph{improbable}, once more
  \item \emph{Say it:} a sentence of your own with \emph{improbable}
  \item \emph{Recall:} say the French for the risk, then the French for risky, then say \emph{improbable} again
\end{itemize}

\begin{tcolorbox}[breakable,colback=teal!4,colframe=teal!35!black,title={Before you move on}]
What does improbable mean? (\textquotedblleft{}Improbable, unlikely\textquotedblright{}.) Say it once more.
\end{tcolorbox}
\section[la probabilité]{la probabilité — the probability, the likelihood}

\label{lesson:FR-C372-probabilite}



\begin{quote}
Before the new one: say the French for risky, then the French for improbable.
\end{quote}

\subsection*{You'll want to know: la probabilité}
\textbf{la probabilité} — \textquotedblleft{}the probability, the likelihood\textquotedblright{}.

\begin{quote}
La probabilité de gagner est très faible.
\end{quote}

\begin{quote}
\emph{The probability of winning is very low.}
\end{quote}

\emph{Selon toute probabilité} — in all likelihood.

\begin{cousinweb}
\emph{probabilité} is a learned borrowing from Latin \emph{probābilitās} \textquotedblleft{}likelihood\textquotedblright{} (Lewis \& Short), from \emph{probābilis} \textquotedblleft{}likely\textquotedblright{}, built on \emph{probāre} \textquotedblleft{}to test\textquotedblright{}.
\end{cousinweb}

\begin{grammarlens}[title={What you've built}]
One more word for telling, explaining and describing at B1.
\end{grammarlens}

\subsection*{Your turn}
\begin{itemize}
\raggedright
  \item \emph{Say it:} \emph{la probabilité}
  \item \emph{Say it:} \emph{la probabilité}, once more
  \item \emph{Say it:} a sentence of your own with \emph{la probabilité}
  \item \emph{Recall:} say the French for risky, then the French for improbable, then say \emph{la probabilité} again
\end{itemize}

\begin{tcolorbox}[breakable,colback=teal!4,colframe=teal!35!black,title={Before you move on}]
What does la probabilité mean? (\textquotedblleft{}The probability, the likelihood\textquotedblright{}.) Say it once more.
\end{tcolorbox}
\section[quitte à]{quitte à — even if it means}

\label{lesson:FR-C372-quitte-a}



\begin{quote}
Before the new one: say the French for improbable, then the French for the probability.
\end{quote}

\subsection*{You'll want to know: quitte à}
\textbf{quitte à} — \textquotedblleft{}even if it means\textquotedblright{}.

\begin{quote}
Je pars tôt, quitte à attendre une heure sur le quai.
\end{quote}

\begin{quote}
\emph{I am leaving early, even if it means waiting an hour on the platform.}
\end{quote}

+ infinitive: you accept the cost. \emph{Être quitte} — to be even, owe nothing.

\begin{cousinweb}
\emph{quitte} comes from Old French \emph{quitte} \textquotedblleft{}free, released\textquotedblright{}, taken from Latin \emph{quiētus} \textquotedblleft{}at rest\textquotedblright{} (EtymDB), from \emph{quiēs} \textquotedblleft{}rest\textquotedblright{} (Lewis \& Short); a Medieval Latin form \emph{quittus} is usually placed in between.
\end{cousinweb}

\begin{grammarlens}[title={What you've built}]
One more word for telling, explaining and describing at B1.
\end{grammarlens}

\subsection*{Your turn}
\begin{itemize}
\raggedright
  \item \emph{Say it:} \emph{quitte à}
  \item \emph{Say it:} \emph{quitte à}, once more
  \item \emph{Say it:} a sentence of your own with \emph{quitte à}
  \item \emph{Recall:} say the French for improbable, then the French for the probability, then say \emph{quitte à} again
\end{itemize}

\begin{tcolorbox}[breakable,colback=teal!4,colframe=teal!35!black,title={Before you move on}]
What does quitte à mean? (\textquotedblleft{}Even if it means\textquotedblright{}.) Say it once more.
\end{tcolorbox}
